\chapter{Benchmark Code and Project Statistics}
\label{appendix:benchmarks}

This appendix provides the benchmark code used for empirical evaluation and statistics about the implementation.

\section{Project Statistics}
\label{sec:project-stats}

Table~\ref{tab:loc} summarises the lines of code for each component of the implementation.

\begin{table}[H]
\centering
\begin{tabular}{|l|l|r|r|}
\hline
\textbf{Chapter} & \textbf{Component} & \textbf{Language} & \textbf{LOC} \\
\hline
\multicolumn{4}{|c|}{\textit{Chapter 4: Extensible Array}} \\
\hline
& ResizableBitVector.java & Java & 89 \\
& DynamicArray.java & Java & 156 \\
& ClassicEliasFano.java & Java & 234 \\
& EliasFanoUniverseKnown.java & Java & 312 \\
& EliasFanoBothUnknown.java & Java & 287 \\
& ExtensibleArrayBenchmark.java & Java & 258 \\
& BitvectorBenchmark.java & Java & 247 \\
\hline
\multicolumn{3}{|r|}{\textit{Subtotal (Chapter 4)}} & 1,583 \\
\hline
\multicolumn{4}{|c|}{\textit{Chapter 5: Fast Predecessor Search}} \\
\hline
& All.cpp & C++ & 512 \\
& BranchingBits.java & Java & 187 \\
& classic\_predecessors.hpp & C++ & 423 \\
& run.sh & Shell & 45 \\
\hline
\multicolumn{3}{|r|}{\textit{Subtotal (Chapter 5)}} & 1,167 \\
\hline
\multicolumn{4}{|c|}{\textit{Chapter 6: Dynamic Prefix Sum}} \\
\hline
& PrefixSum.cpp & C++ & 234 \\
& DynamicPrefixSum.hpp & C++ & 189 \\
\hline
\multicolumn{3}{|r|}{\textit{Subtotal (Chapter 6)}} & 423 \\
\hline
\multicolumn{4}{|c|}{\textit{Supporting Code}} \\
\hline
& Build scripts \& configuration & Various & 156 \\
& Test harnesses & Java/C++ & 312 \\
\hline
\multicolumn{3}{|r|}{\textit{Subtotal (Support)}} & 468 \\
\hline
\hline
\multicolumn{3}{|r|}{\textbf{Total}} & \textbf{3,641} \\
\hline
\end{tabular}
\caption{Lines of code by component. Excludes blank lines and comments. Third-party libraries (sux4j, fastutil) are not counted.}
\label{tab:loc}
\end{table}

\section{Extensible Array Benchmark}
\label{sec:ea-benchmark}

\begin{lstlisting}[language=Java, caption={ExtensibleArrayBenchmark.java - Comparing resize strategies}, label={lst:ea-bench}]
package com.chp3.efano;

import it.unimi.dsi.bits.LongArrayBitVector;
import java.util.ArrayList;
import java.util.Random;

public class ExtensibleArrayBenchmark {
    private static final int WARMUP = 5;
    private static final int ITERATIONS = 10;
    private static final int[] SIZES =
        {1000, 10000, 100000, 1000000, 10000000};

    public static void main(String[] args) {
        System.out.println("Size\tRBV(ms)\tArrayList(ms)" +
                          "\tSux4j(ms)\tRBV_Space\tAL_Space");

        for (int size : SIZES) {
            runBenchmark(size);
        }
    }

    private static void runBenchmark(int size) {
        // Warmup
        for (int i = 0; i < WARMUP; i++) {
            benchmarkRBV(size);
            benchmarkArrayList(size);
            benchmarkSux4j(size);
        }

        // Benchmark
        long rbvTime = 0, alTime = 0, sux4jTime = 0;
        for (int i = 0; i < ITERATIONS; i++) {
            rbvTime += benchmarkRBV(size);
            alTime += benchmarkArrayList(size);
            sux4jTime += benchmarkSux4j(size);
        }

        // Space measurement
        long rbvSpace = measureRBVSpace(size);
        long alSpace = measureALSpace(size);

        System.out.printf("%d\t%.2f\t%.2f\t%.2f\t%d\t%d%n",
            size,
            rbvTime / (double) ITERATIONS,
            alTime / (double) ITERATIONS,
            sux4jTime / (double) ITERATIONS,
            rbvSpace, alSpace);
    }

    private static long benchmarkRBV(int size) {
        ResizableBitVector rbv = new ResizableBitVector();
        Random rand = new Random(42);
        long start = System.nanoTime();
        for (int i = 0; i < size; i++) {
            rbv.append(rand.nextInt(2));
        }
        return (System.nanoTime() - start) / 1_000_000;
    }

    private static long benchmarkArrayList(int size) {
        ArrayList<Integer> list = new ArrayList<>();
        Random rand = new Random(42);
        long start = System.nanoTime();
        for (int i = 0; i < size; i++) {
            list.add(rand.nextInt(2));
        }
        return (System.nanoTime() - start) / 1_000_000;
    }

    private static long benchmarkSux4j(int size) {
        LongArrayBitVector bv =
            LongArrayBitVector.getInstance();
        Random rand = new Random(42);
        long start = System.nanoTime();
        for (int i = 0; i < size; i++) {
            bv.add(rand.nextInt(2) == 1);
        }
        return (System.nanoTime() - start) / 1_000_000;
    }
}
\end{lstlisting}

\section{Predecessor Search Benchmark}
\label{sec:pred-benchmark}

\begin{lstlisting}[language=C++, caption={Predecessor search timing benchmark}, label={lst:pred-bench}]
#include <chrono>
#include <random>
#include <iostream>
#include <immintrin.h>

// Number of queries for timing
const int NUM_QUERIES = 1 << 27;

// Benchmark function template
template<typename Func>
double benchmark(Func f, const uint16_t* S, int k) {
    std::mt19937 gen(42);
    std::uniform_int_distribution<uint16_t> dist(0, 65535);

    auto start = std::chrono::high_resolution_clock::now();

    volatile long sum = 0;  // Prevent optimisation
    for (int i = 0; i < NUM_QUERIES; i++) {
        uint16_t x = dist(gen);
        sum += f(S, x);
    }

    auto end = std::chrono::high_resolution_clock::now();
    auto duration =
        std::chrono::duration_cast<std::chrono::microseconds>
        (end - start);

    return duration.count() / 1e6;  // Return seconds
}

int main(int argc, char* argv[]) {
    int k = (argc > 1) ? atoi(argv[1]) : 16;

    // Generate random sorted set
    std::vector<uint16_t> S(k);
    std::mt19937 gen(123);
    std::uniform_int_distribution<uint16_t> dist(0, 65535);
    for (int i = 0; i < k; i++) {
        S[i] = dist(gen);
    }
    std::sort(S.begin(), S.end());

    // Run benchmarks
    std::cout << "k\ttLinear\ttBSI\ttAFK\ttFWM\ttSTL\n";
    std::cout << k << "\t";

    // Linear search
    std::cout << benchmark(rankLinear<16>, S.data(), k);
    std::cout << "\t";

    // Bit-slice index
    std::cout << benchmark(rankBSI<16>, S.data(), k);
    std::cout << "\t";

    // Ajtai with lookup tables
    std::cout << benchmark(rankAFK<16>, S.data(), k);
    std::cout << "\t";

    // Improved Fredman with XOR-MinPos
    std::cout << benchmark(rankFWImproved<16>, S.data(), k);
    std::cout << "\t";

    // STL lower_bound
    std::cout << benchmark(rankSTL<16>, S.data(), k);
    std::cout << "\n";

    return 0;
}
\end{lstlisting}

\section{Build and Run Instructions}
\label{sec:build}

\subsection{Chapter 4: Elias-Fano (Java/Maven)}

\begin{lstlisting}[language=bash, caption={Building and running Chapter 4 experiments}]
cd final_thesis_experiments-main/chp3/EliasFano

# Compile
mvn compile

# Run benchmarks
mvn exec:java -Dexec.mainClass=\
  "com.chp3.efano.ExtensibleArrayBenchmark"

# Run bitvector benchmarks
mvn exec:java -Dexec.mainClass=\
  "com.chp3.efano.BitvectorBenchmark"
\end{lstlisting}

\subsection{Chapter 5: Fast Predecessor Search (C++)}

\begin{lstlisting}[language=bash, caption={Building and running Chapter 5 experiments}]
cd final_thesis_experiments-main/chp4/FastInSmall

# Compile with AVX2 optimisations
g++ -std=c++11 -march=native -mssse3 -mavx2 \
    -mbmi2 -O3 All.cpp -o Demo

# Run for specific set size
./Demo 16

# Run full benchmark suite
./run.sh
\end{lstlisting}

\section{Computing Environment}
\label{sec:environment}

All experiments were conducted on the following system:

\begin{table}[H]
\centering
\begin{tabular}{|l|l|}
\hline
\textbf{Component} & \textbf{Specification} \\
\hline
CPU & Intel Core i5-4460 @ 3.40GHz \\
Architecture & x86\_64 (64-bit) \\
SIMD Support & SSE4.2, AVX, AVX2, BMI1, BMI2 \\
RAM & 8 GB DDR3 \\
L1 Cache & 32 KB (data), 32 KB (instruction) \\
L2 Cache & 256 KB \\
L3 Cache & 6 MB \\
\hline
OS & Ubuntu 22.04.3 LTS \\
Kernel & Linux 5.15.0 \\
Java & OpenJDK 11.0.20 \\
C++ Compiler & g++ 11.4.0 \\
Optimisation & -O3 -march=native \\
\hline
\end{tabular}
\caption{Computing environment for experimental evaluation}
\label{tab:environment}
\end{table}

\section{Data Availability}
\label{sec:data}

The complete source code and benchmark data are available at:

\begin{itemize}
    \item \textbf{Repository}: The implementation code is organised in the \texttt{final\_thesis\_experiments-main} directory.

    \item \textbf{Benchmark Results}: Raw benchmark data in TSV format is stored in \texttt{result.tsv} files within each chapter's directory.

    \item \textbf{Reproducibility}: All experiments can be reproduced by running the provided shell scripts (\texttt{run.sh}) after installing the required dependencies.
\end{itemize}
